\pagestyle{plain} % No headers, just page numbers
\pagenumbering{arabic} % Arabic numerals
\setcounter{page}{1}

\chapter{\uppercase {Introduction}}

A wide range of computational problems in science and engineering naturally
arise as combinatorial optimization (CO) problems on graphs. From
identifying influential individuals in social
networks~\cite{kempe2003maximizing} to designing efficient transportation
routes~\cite{sun2021generalization} and allocating resources in telecommunication
systems, graph CO problems are pervasive. These problems are, in general, computationally intractable: among the 21 NP-complete
problems introduced by Karp~\cite{karp2009reducibility}, 10 are decision versions of
graph optimization problems, and most of the remaining ones can be naturally reformulated
on graphs. Over
the decades, researchers have developed heuristics and approximation algorithms that
produce near-optimal solutions in polynomial
time~\cite{nemhauser1978analysis, kempe2003maximizing, tang2015influence}, enabling
practical solution on problems that are intractable to solve exactly.

As real-world graphs continue to grow in size, even state-of-the-art solvers face
significant scalability challenges. Modern real-world graphs routinely contain
millions or billions of nodes, from online social networks to web-scale knowledge
graphs. At the same time, the budget constraints of many practical applications are such
that an optimal solution involves only a small number of elements. For example, in viral
marketing, a company may have the budget to seed only a handful of individuals while
aiming to maximize overall information spread~\cite{kempe2003maximizing}. In these settings,
the vast majority of candidates in the ground set are irrelevant to the optimal
solution, yet existing solvers must consider all of them. This mismatch between the
enormous size of the ground set and the small size of the optimal solution creates a
fundamental scalability bottleneck. A natural strategy to address this bottleneck is pruning: given the full ground set
$\mathcal{U}$, the goal is to identify and discard elements that are unlikely to
contribute to a high-quality solution, producing a much smaller subset 
$\mathcal{U}' \subseteq \mathcal{U}$ with $|\mathcal{U}'| \ll |\mathcal{U}|$. A solver can then operate on $\mathcal{U}'$ instead of $\mathcal{U}$, ideally with
little or no loss in solution quality. Formally, a CO problem is a tuple
$(S, \Omega, f)$, where $S$ is the search space over discrete decision variables,
$\Omega$ the feasibility constraints, and $f: S \to \mathbb{R}$ the objective function.
Given a ground set $\mathcal{U}$ and a solver $\mathcal{OPT}$, the optimal solution is
%
\begin{equation}
  X^\star = \mathcal{OPT}(\mathcal{U}) = \arg\max_{X \subseteq \mathcal{U}, \; X \text{ satisfies } \Omega} f(X).
\end{equation}
%
The pruning problem seeks a subset $\mathcal{U}' \subseteq \mathcal{U}$ with
$|\mathcal{U}'| \ll |\mathcal{U}|$ such that
%
\begin{equation}
  \frac{f\bigl(\mathcal{OPT}(\mathcal{U}')\bigr)}{f\bigl(\mathcal{OPT}(\mathcal{U})\bigr)} \geq 1 - \epsilon,
\end{equation}
%
for a small approximation factor $\epsilon \geq 0$. However, existing
pruning approaches are limited by poor scalability of sequential classical
methods~\cite{zhou2017scaling}, applicability to only a single CO problem in methods designed for maximum clique
enumeration~\cite{lauri2019fine, grassia2019learning, lauri2023learning} or the traveling
salesman problem~\cite{fitzpatrick2021learning, sun2021generalization}, reliance on
hand-crafted, problem-specific features in learning-based
methods~\cite{manchanda2020gcomb, ireland2022lense, tian2024combhelper}, and a complete
absence of theoretical guarantees on pruning quality.
Altogether, these limitations motivate a focused research effort on developing pruning
strategies that are scalable, general, and principled. The overarching goal of the
proposed work is to address this challenge by developing a suite of pruning methods for
graph CO, aiming to make solvers faster, more general, and
more practical for large-scale graph instances.

Prior to our first work, all existing pruning methods were heuristics with no formal
performance guarantees. It was unknown whether it is even possible to efficiently produce
a pruned ground set with provable bounds on both its size and the fraction of optimal
value it retains. To answer this question, we introduce
QuickPrune~\cite{pmlr-v258-nath25a}, a lightweight pruning algorithm for constrained
discrete optimization that processes elements in a single pass through the ground set,
discarding those unlikely to appear in an optimal solution. Under mild assumptions on
the objective function and element costs, we prove that the pruned ground set retains a
constant fraction of the optimal value across a range of budgets, while its size grows
only logarithmically in the original ground set. QuickPrune is the first pruning
algorithm to provide such guarantees, and the first to support a range of budgets
simultaneously rather than a single budget instance. Through extensive experiments on
real-world datasets for various applications, we demonstrate that QuickPrune efficiently
prunes the ground set and outperforms state-of-the-art classical and machine learning
heuristics for pruning.

As graph instances grow to millions of nodes, even a lightweight sequential algorithm
faces scalability challenges, particularly for problems where calculating marginal gain
is \#P-hard. At the same time, existing learning-based pruning methods achieve fast
inference but demand problem-specific domain knowledge for algorithm design, making them
difficult to apply to new CO problems. To overcome both of these limitations, we propose
Hierarchical DeepPruner (H-DeepPruner)~\cite{nath2025hdeeppruner}, an adaptive
two-stage framework that combines supervised learning with reinforcement learning for
search space reduction under size constraints. The first stage uses a graph neural
network to quickly eliminate vertices unlikely to be part of the solution, without
requiring any hand-crafted features or domain knowledge. The second stage applies a
tree search procedure to further refine the pruned ground set by considering
interdependencies among elements. Our experiments demonstrate that H-DeepPruner
significantly outperforms existing learning-based pruning methods across multiple graph
CO problems, showing that principled combinations of supervised and reinforcement
learning can achieve effective pruning without domain-specific knowledge.

While H-DeepPruner eliminates the need for hand-crafted features, the features it uses
are problem-agnostic and may lack the expressiveness needed to capture problem-specific
structural patterns. On the other end of the spectrum, existing methods use carefully
engineered features that are expressive but require substantial domain expertise,
limiting their transferability across problems and constraints. To bridge this gap, we
propose LLM2Prune~\cite{nath2025llm2prune}, a framework that leverages Large Language
Models to automatically discover problem-grounded node features for a downstream
classifier. Given only a natural language task description, LLM2Prune guides the feature
discovery process using both performance evaluation and feature-importance feedback,
producing features that are tailored to the specific problem without any human
engineering. LLM2Prune discovers non-trivial features automatically that would typically
require significant domain expertise to design manually, and adapts naturally to
different constraint types by simply modifying the task description. Our experiments
show that LLM2Prune achieves competitive performance with classical methods while
consistently outperforming existing learning-based pruning approaches.

The research presented in this proposal collectively addresses the challenge of scaling
combinatorial optimization to massive graph instances through principled search space
reduction. QuickPrune establishes the theoretical foundations, proving that pruning with
formal value-retention bounds is achievable under natural assumptions. Hierarchical
DeepPruner demonstrates that learning-based methods can achieve superior empirical
performance under size constraints without requiring any domain expertise. LLM2Prune
completes the picture by automating the feature discovery process entirely, closing the
gap between problem-agnostic design and problem-aware performance. This integrated
perspective has implications for building combinatorial optimization solvers that are
not only faster through effective search space reduction, but also more general across
problems and constraints, and more practical through automated feature discovery.

The remainder of this proposal is organized as follows. Chapter~2 reviews related work
on pruning methods for combinatorial optimization. Chapters~3, 4, and 5 present
QuickPrune, Hierarchical DeepPruner, and LLM2Prune in detail, respectively. Chapter~6
discusses proposed future work and open research directions, and Chapter~7 presents the
research timeline.